\documentclass[11pt]{article}
\usepackage[margin=1in]{geometry}
\usepackage{times}
\usepackage{hyperref}
\hypersetup{colorlinks=true, linkcolor=black, urlcolor=blue, citecolor=black}
\title{Mechanical Power, Electrical Input, and the N-Machine Claim}
\author{Flyxion \\ \small Independent Researcher}
\date{}
\begin{document}
\maketitle

\begin{abstract}
Bruce DePalma's N-machine is the clearest example in this volume of an energy claim that depends on treating a proxy as if it were the quantity required by the power balance. DePalma described the device as a one-piece rotating magnetized conductor that could deliver electrical output with far less back torque than a conventional generator delivering the same current, a claim that, if correct, would place it in the over-unity category. This paper follows the actual historical record of the claim, from DePalma's own descriptions of the Sunburst machine in Santa Barbara in the late 1970s to Kincheloe's mid-1980s examination of the device, and explains why the central issue is not simply whether the machine was unusual, but whether it was measured at the shaft. The answer is that the published record does not contain the kind of independent, calibrated torque-sensor test that would justify a claim of over-unity, and the well-established question about unipolar induction does not change the energy accounting at the shaft. The paper therefore treats the N-machine as a claim whose apparent power gain depends on a specific set of measurement practices, not as a confirmed violation of the conservation of energy.
\end{abstract}

\section{Introduction}

The N-machine belongs to a family of claims in which a highly specific rotating electromagnetic geometry is treated as if it automatically overturns the ordinary law of power conservation. In the DePalma case, the claim was advanced as a machine in which a one-piece rotating magnetized conductor could extract electrical power with far less mechanical opposition than a conventional generator would show under the same current. The practical relevance of the claim, and the reason it keeps resurfacing in the literature, is that the power balance in a rotating electric machine is not judged by the electrical input alone. It is judged by the torque and angular speed at the shaft. If the machine is described as having unusually low back torque, the claim must be supported by a direct mechanical measurement at the shaft. Otherwise one is not measuring the machine's actual conversion efficiency but a convenience quantity that can be optimistic for reasons that have nothing to do with the real rotational load.

Bruce DePalma was an American electrical engineer and inventor best known for the N-machine, a rotating electrical machine he claimed could extract additional energy from space and produce more electrical power than the mechanical power supplied to drive it. The device was a homopolar generator in the conventional sense: a conducting disc rotating in a magnetic field develops a voltage between its centre and rim. The established physics of that arrangement is not in dispute. The disputed point is whether the specific machine geometry and brush arrangement could deliver a net surplus once the true shaft power, losses, and back torque were measured. That is why the N-machine belongs to the broader family of over-unity claims in which the central problem is not the basic electromechanics but the power accounting.

The historical record is therefore more important than the rhetoric. DePalma's published descriptions of the N-machine are tied to his broader program of rotating-field and unidirectional-force claims \cite{DP1974,DP1990}. The machine that most clearly carries that program into the electromechanical domain is the Sunburst machine, a one-piece rotating magnet design associated with DePalma in Santa Barbara around 1978 and 1979. DePalma reported that it showed far less back torque than a conventional generator delivering the same current, roughly thirteen to twenty percent of the expected drag \cite{DP1990}. That is an extraordinary claim, and it is the claim this paper assesses. The question is not whether a rotating magnetized conductor raises a real conceptual problem about how flux is to be regarded in a one-piece rotor. It does. The question is whether that conceptual issue changes the mechanical work required to deliver current to the load. It does not, because the power required to sustain the rotor remains a function of torque and angular speed at the shaft.

The pattern the N-machine exposes is not unique to DePalma. It appears wherever an inventor measures the drive motor's electrical input, or the source power feeding the apparatus, and treats that as if it were the real shaft power. Once the measurement is done at the wrong place in the system, a machine can appear to outperform its own mechanical input even though the real shaft torque has not been measured. That is why the most important evidence in the N-machine story is not the dramatic claim itself but the record of who actually measured what, and with what instrument.

\section{History of the Claim}

Bruce DePalma's N-machine was not an isolated oddity. It belonged to a longer design program developed in the context of his broader interest in the electromagnetic and mechanical implications of rotation. His 1974 monograph, \emph{The Generation of a Unidirectional Force}, set out the general conceptual framework in which rotation, magnetization, and force generation were treated as physically linked in a way that standard electromechanical theory did not accept \cite{DP1974}. The monograph is not a conventional engineering design document, and it is not the sort of work that can be evaluated as a mature machine specification. It is a conceptual statement of a program. The N-machine extended that program into a more concrete apparatus: a rotating, magnetized conductor in which the rotor itself served as part of the active electrical structure rather than as a separate armature and field system.

The central design feature was deliberately simple. The rotor was not built as a conventional field winding plus armature assembly. It was a one-piece rotating conductor with a magnetized structure or magnetized rotor. According to DePalma's own descriptions, the objective was to reduce the back torque that conventional generator systems impose on the shaft when current is drawn from the machine. DePalma stated that the Sunburst machine, built in Santa Barbara around 1978 to 1979, showed far less back torque than a conventional generator delivering the same current, amounting to roughly thirteen to twenty percent of the expected drag \cite{DP1990}. That was not a small engineering improvement. It was a claim that the apparatus was producing useful power with a mechanical opposition far below the ordinary magnetic braking expected in a generator of comparable output.

That specific claim is the relevant one. The historical context matters in part because it explains why the N-machine attracted attention. It arrived amid a broader movement in which unconventional rotating-electrical claims were being treated as possible evidence of hidden energy flows. But the historical context does not remove the requirement for shaft-power measurement. The design premise may be unusual, but the power balance remains ordinary. If the current is being drawn from the machine, the source of the electrical energy must still be the mechanical work delivered to the shaft, with losses accounted for at the rotor, brushes, bearing, drive, and any internal magnetic drag. A claim of unusually low back torque is therefore not evidence of free energy; it is only a sign that a particular machine may have lower parasitic drag than a conventional generator. The difference between a low-loss machine and a free-energy machine is not a rhetorical distinction. It is a measurement distinction.

The question of whether the one-piece rotating magnetized conductor raised a deeper issue in electromagnetic theory also matters, and it should be stated honestly. Unipolar induction, the old problem of whether the magnetic flux is considered to rotate with the conductor or to remain fixed in space, is a genuine conceptual puzzle in electromagnetism. It is not a trivial matter. It is a longstanding problem in rotating-frame electrodynamics, and it is the reason that physically unusual rotor arrangements can invite a closer look at where the electromagnetic fields are actually moving relative to the conducting material \cite{Jackson1999}. That issue is real. But it is not an argument for excess power. It is an argument about how best to describe rotating electromagnetic fields in a frame-dependent formalism. A real conceptual question about the field's motion does not, by itself, alter the mechanical work required to keep the rotor turning against the generated current.

\section{The Actual Test Record}

The strongest evidence in the N-machine saga comes from the test record rather than from the inventor's descriptive rhetoric. The crucial figure is Robert Kincheloe, whose examination of the Sunburst machine in the mid-1980s was reported in the DePalma literature and later discussed in the pages of \emph{The Skeptical Inquirer} in volume 15 \cite{K1986,SI1991}. Kincheloe's central observation was that DePalma's 1980 power measurements on the alternating-current drive did not account for phase angle as the load changed. That matters because in an alternating-current circuit, power is not simply voltage times current. It is the product of voltage, current, and the cosine of the phase angle. If the phase angle changes with load, then the power calculation changes as well, and a comparison that ignores the angle can become misleadingly favorable to the machine under test. The point is not merely that the measurement was imprecise; it is that an a.c. power calculation is not valid unless the phase angle is included, and DePalma's reported comparison effectively treated motor input as if it were shaft power without that correction.

This is not a minor detail. It is exactly the kind of measurement issue that can produce an optimistic result when a machine is evaluated by a simplified electrical reading. A current reading and a voltage reading do not determine real power unless one includes phase. If the motor load or the generator output shifts phase as the system changes, then a comparison between the motor input and the apparent output that ignores phase is not an energy balance. It is an estimate that can easily show an apparent surplus. The point made by Kincheloe was therefore not that the machine was a proven free-energy device, but that the published numbers did not satisfy the minimum instrumentation requirement needed to support that claim. The claim remained, at best, untested under the standard it required.

The same point applies to later works in the same family. The later large N-1 machine, discussed in the broader DePalma literature, did not establish the sort of independent shaft-power test that would resolve the question. Nor did the later Tewari work supply the independent mechanical torque record that would make the power claim credible. In the Tewari case, the relevant issue is the same one that appears across several over-unity proposals: a claim of unusual electrical output is made without a measured torque-and-speed balance at the shaft, and the published documentation does not show that standard machinery design checks were applied to the actual mechanical input. That is not proof that the machine is wrong. It is evidence that the claim remains untested at the standard the claim requires.

This is an important distinction. The absence of a published, controlled torque-sensor test of the N-machine does not establish that the machine is impossible. It establishes that the machine's power claim remains unresolved under the measurement standard required to judge it. The correct phrasing is therefore narrow and exact: no independent, published shaft-power test is in the record I can verify, and so the claim remains unsubstantiated. That is not a refutation. It is a statement of evidentiary status.

\section{The Physics Fairly Stated}

A one-piece rotating magnetized conductor does raise a real question about where the electromagnetic field is considered to rotate and where the current is generated. The older unipolar induction issue is exactly the kind of problem that makes a geometry like this look strange to someone who expects the conventional armature-field separation to be universal. If the magnet rotates with the conductor, is the field rotating with it, or is it fixed in a laboratory frame? The answer depends on the chosen reference frame and the path one uses to compute the induced electromotive force. That is a genuine theoretical issue, and it deserves to be presented honestly rather than dismissed as a trivial quibble. The history of electrical engineering is full of apparently simple rotating conductor problems that become conceptual puzzles when one tries to insist that the field must be attached to one frame or another.

But this conceptual issue does not alter the existence of the shaft-power requirement. A machine can show an unusual field pattern and still obey the ordinary relation between electrical output and mechanical input. The force required to drive the rotor is the mechanical torque times the angular speed. The generated electrical power is the load power delivered. The back torque is a mechanical drag that must be supplied by the drive. If the field arrangement lowers the magnetic drag, the machine is simply lower-loss; it is not free-energy. A rotation of the field is not the same thing as the creation of power. The bookkeeping at the shaft remains the governing constraint. The conceptual complication is real, but it changes the description of the electromagnetic geometry, not the law governing the mechanical-to-electrical conversion.

That is the crucial distinction. A one-piece rotating conductor may look like a challenge to the ordinary generator picture, and it may deserve a careful field analysis. But a field analysis is not a power claim. The power claim still has to survive a shaft power test, measured torque and angular velocity, with all losses included. That is the standard the N-machine claim did not meet in the published record.

\section{The Entropic Gravity Framing}

The N-machine is an energy claim, not a gravitational one, and the entropic-gravity framing used elsewhere in this series is therefore intentionally brief. In an entropic account of gravitation, spacetime curvature is treated as the large-scale expression of an ordered configuration, not as a field that can be shielded or filtered by a rotating conductor in the way a magnetic field can. A rotating magnetized conductor can alter the electrical and mechanical conditions of a machine, but it does not generate a new gravitational mechanism by virtue of rotating. The issue at stake here is therefore not the structure of gravitation but the measurement of electromechanical power. The energy-accounting problem remains the operative one, and the shaft-power requirement is the decisive test.

\section{Objections}

The strongest defender of the N-machine argument would say that DePalma's design was not a conventional generator at all but a genuinely new electromechanical architecture, and that a one-piece rotating magnetized conductor could yield a lower back torque than ordinary machines because the flux was not moving in the same way as it does in a wound armature. That is a coherent objection. It is also exactly the kind of objection that must be tested with shaft-power instrumentation rather than with simplified motor-input readings. An unusual field geometry is not evidence of anomalous power unless the shaft torque and speed are measured and the full loss budget is accounted for.

A second objection is that the absence of a published controlled torque-sensor test is not the same as a failed test. That is correct, but it is not the same as a successful test either. The correct conclusion is that the claim remains untested at the standard it requires. That is the strongest statement the record supports. It is also the one that keeps the question open without turning unsupported numbers into evidence.

\section{Conclusion}

The N-machine claim depends on a specific measurement mistake: the use of a convenient electrical proxy in place of the actual machine input, which is torque at the shaft times angular speed. DePalma's Sunburst machine was described as showing only roughly thirteen to twenty percent of the expected back torque, but the published historical record does not show the independent shaft-power test that would justify the extraordinary conclusion. Kincheloe's mid-1980s examination of the Sunburst machine identified the missing phase-angle correction in the alternating-current drive, and that observation matters because in an a.c. circuit the power calculation cannot be made without phase angle. The later large N-1 machine and the later Tewari work are best read in the same way: claims exist, but the independent, calibrated shaft-power measurements needed to establish them are not part of the public record I have been able to verify. This paper therefore treats the N-machine as a historically significant but not yet validated claim, one whose mechanical question is real, but whose energy claim remains unsupported by the standard of proof it requires.

\begin{thebibliography}{99}

\bibitem{DP1974} B. E. DePalma, \emph{The Generation of a Unidirectional Force}, Simularity Institute, 1974.

\bibitem{DP1990} B. DePalma, On the Possibility of Extraction of Electrical Energy Directly from Space, \emph{Speculations in Science and Technology} 13 (1990) 283.

\bibitem{K1986} R. Kincheloe, Homopolar ``Free-Energy'' Generator Test, report on the Sunburst N machine, as discussed in the DePalma literature and later summarized in \emph{The Skeptical Inquirer} 15 (1990--1991) 72.

\bibitem{SI1991} \emph{The Skeptical Inquirer} 15 (1990--1991) 72.

\bibitem{Jackson1999} J. D. Jackson, \emph{Classical Electrodynamics}, 3rd ed., Wiley, 1999.

\end{thebibliography}

\end{document}
